%\section{Introduction}
\section{Background} 
Skin conditions affect people of different ages and can influence both physical health and quality of life. Conditions such as acne, eczema, and burns may vary in their appearance and severity, making accurate identification important for appropriate care. However, access to dermatological services can be difficult in some areas because of the limited availability of specialists, consultation costs, and geographical barriers. These challenges have encouraged researchers to investigate computer-aided systems that can analyze skin images and assist with the identification of skin conditions. Such systems may provide useful preliminary information, although they cannot replace professional medical diagnosis.

Recent developments in artificial intelligence have created new opportunities for analyzing medical images. Convolutional Neural Networks (CNNs) can learn visual features from images and use them to distinguish between different categories. Transfer learning provides another approach by allowing models trained on large image datasets to be adapted to specific classification tasks. Models such as EfficientNetB0, MobileNetV2, and ResNet50 offer different architectural approaches to feature extraction and computational efficiency. However, their performance may vary depending on the dataset, the classification task, and the proportions of data used for training, validation, and testing. Therefore, comparing these models under different data-splitting configurations can help identify a suitable approach for skin condition and severity classification.

Although deep learning models can classify images, their predictions do not always make it clear which visual features influenced the results. This limitation is particularly important in medical image analysis, where understanding a model's predictions is necessary for evaluating its usefulness. Explainable Artificial Intelligence (XAI) methods can help address this issue by providing visual information about a model's decision-making process. One such method is Gradient-weighted Class Activation Mapping (Grad-CAM), which produces a heatmap highlighting image regions that contribute to a selected model prediction. In this study, Grad-CAM is used to examine the predictions made by the evaluated models, alongside a comparison of their classification performance and training time. By investigating different model architectures and dataset splits, this research aims to assess their suitability for skin condition classification and severity identification while providing visual explanations of the predictions.

\section{Problem Statement}
Conventional diagnosis of skin conditions primarily depends on visual examination by experienced dermatologists. Although this approach is effective, it may not always be readily available, especially in remote or resource-limited areas. Delays in diagnosis can result in disease progression and reduced treatment effectiveness. In addition, many existing automated skin analysis systems focus solely on disease classification without providing information about severity or explaining how predictions are made.\\
Another important limitation is the lack of practical tools for monitoring changes in skin conditions over time. Patients undergoing treatment often rely on subjective visual observation to determine whether their condition is improving. An intelligent system capable of estimating severity levels from skin images can provide a consistent and objective method for comparing images captured at different stages of treatment. Such a system can assist healthcare professionals and patients in monitoring disease progression or recovery while improving confidence through visual explanations.

\section{Research Objectives}
The primary objective of this research is to develop an explainable deep learning framework for automatic skin condition classification and severity estimation using transfer learning techniques. \\

The specific objectives of this study are:
\begin{itemize}
    \item To develop an automated skin condition classification system capable of identifying multiple skin conditions from image data.
    \item To estimate the severity or condition level associated with each identified skin condition.
    \item To implement and compare multiple transfer learning models, including EfficientNetB0, MobileNetV2, and ResNet50, for skin condition classification.
    \item To generate confidence scores for every skin condition–severity label, providing a comprehensive prediction rather than a single output.
    \item To incorporate Grad-CAM visualization to explain the model's decision-making process by highlighting influential image regions.
\end{itemize}
To provide a framework that can assist in monitoring changes in skin condition severity over time by comparing predictions obtained from images captured during different stages of treatment.

\section{Scope of the Study}
This research focuses on the classification of four skin condition categories: Acne, Burn, Eczema, and Normal skin. The study utilizes an image dataset with multiple severity or condition labels and employs transfer learning to train deep learning models. Three pre-trained CNN architectures - EfficientNetB0, MobileNetV2, and ResNet50 - are comparatively evaluated to determine their effectiveness for the proposed task.\\
The developed framework predicts confidence scores for each skin condition - severity label and provides Grad-CAM visualizations to improve model interpretability. The proposed system is intended to function as a computer-aided decision support tool and is not designed to replace professional medical diagnosis. Instead, it aims to assist users by providing preliminary assessments and visual explanations while supporting the monitoring of changes in skin condition severity throughout the treatment process.

\section{Organization of the Thesis}
The remainder of this thesis is organized as follows. Chapter 2 presents a review of the relevant literature on skin condition classification, transfer learning, and explainable artificial intelligence. Chapter 3 describes the research methodology, including dataset preparation, preprocessing, model development, training procedures, and evaluation metrics. Chapter 4 presents the experimental results, comparative analysis of the selected transfer learning models, and Grad-CAM visualizations. Finally, Chapter 5 concludes the study by summarizing the major findings, discussing limitations, and providing recommendations for future research.


% \nomenclature{$ODI$}{One day International}
% \nomenclature{$T20$}{Twenty Twenty}
% \nomenclature{$ODI$}{One day International}
% \nomenclature{$IPL$}{Indian Premier League}
% \nomenclature{$ODI$}{One day International}
% \nomenclature{$MR$}{Runs scored by Home team}
% \nomenclature{$OR$}{Runs scored by the opponent team}
% \nomenclature{$MRN$}{Home Team Run Rate}
% \nomenclature{$ORN$}{Opponent Team Run Rate}
% \nomenclature{$LBW$}{Leg before Wicket}